%% A descending construction with its exact choice assumptions.

\section{The Rasiowa--Sikorski dense-set lemma}

\begin{annexitem}{annex:p10.rs.1}{thm:p10-rasiowa-sikorski}{A generic filter meeting countably many dense sets}
\begin{proof}
\paragraph{The descending sequence under dependent choice.}
First suppose \(\mathcal D\) is presented as
\(\{D_n:n\in\mathbb N\}\), allowing repetitions.
On the nonempty set \(S=\mathbb N\times P\), define
\[
 (n,q)R(m,r)
 \quad\Longleftrightarrow\quad
 m=n+1,\quad r\leq q,\quad r\in D_n.
\]
For every \((n,q)\), density of \(D_n\) supplies such an \(r\),
so the relation is serial.
Apply the stated form of \(\mathsf{DC}\) with prescribed start
\((0,p)\).
Induction on the first coordinates shows that the resulting sequence
has the form \((n,p_n)\), where
\[
 p_0=p,\qquad p_{n+1}\leq p_n,\qquad p_{n+1}\in D_n.
\]

\paragraph{The upward closure is the required filter.}
Put
\[
 G=\{q\in P:\exists n\quad p_n\leq q\}.
\]
It contains \(p_0=p\), hence is nonempty.
If \(q\in G\) and \(q\leq r\), a witness \(p_n\leq q\)
also satisfies \(p_n\leq r\), so \(r\in G\).
For \(q,r\in G\), take witnesses \(p_i\leq q\) and
\(p_j\leq r\), and set \(k=\max(i,j)\).
The descending chain gives
\(p_k\leq p_i\leq q\) and \(p_k\leq p_j\leq r\).
Since \(p_k\in G\), this proves downward directedness.
Finally, \(p_{n+1}\) belongs to both \(G\) and \(D_n\).
Thus \(G\) is \(\mathcal D\)-generic.

\paragraph{A fixed enumeration of the poset.}
Suppose \(e:\mathbb N\to P\) is onto.
For every \(n\) and \(q\in P\), the set
\[
 A_{n,q}=\{j\in\mathbb N:e(j)\leq q\text{ and }e(j)\in D_n\}
\]
is nonempty, by density and surjectivity of \(e\).
Define uniquely by recursion in \(\mathsf{ZF}\)
\[
 p_0=p,\qquad p_{n+1}=e(\min A_{n,p_n}).
\]
This sequence has the same descent and membership properties, so
the preceding filter verification applies without \(\mathsf{DC}\).
Taking a least index specifies one value at every step; no sequence
of unspecified choices is used.  This is a set-theoretic construction
and does not require the order or the dense sets to be computably
decidable.

\paragraph{Finite and empty families.}
If \(\mathcal D=\{D_0,\ldots,D_{m-1}\}\), density and finite
induction give
\[
 p_0=p,\qquad p_{i+1}\leq p_i,\quad p_{i+1}\in D_i
                   \qquad(0\leq i<m).
\]
Only finitely many choices are required, which is provable in
\(\mathsf{ZF}\).
The principal upward set
\[
 G=\{q\in P:p_m\leq q\}
\]
is a filter containing \(p\): its common stronger condition for
any two members is \(p_m\).
For each \(i<m\), the earlier witness \(p_{i+1}\) lies in
\(D_i\) and in \(G\), since \(p_m\leq p_{i+1}\).
Thus downward closure of the dense sets is unnecessary.
When \(m=0\), take \(G=\{q:p\leq q\}\); genericity is vacuous.

Every nonempty countable family has an enumeration, with repetitions
permitted in the finite case. Choosing one such enumeration suffices
to apply the sequence argument. The empty family has just been
handled separately.

\paragraph{Finite posets.}
If \(P\) is finite, the nonempty finite set
\(\{q\in P:q\leq p\}\) has a minimal member \(r\).
Every \(s\leq r\) also lies below \(p\), so minimality implies
\(s=r\).
For any downward dense \(D\subseteq P\), density below \(r\)
therefore gives \(r\in D\).
The filter \(\{q:r\leq q\}\) contains \(p\) and meets every
dense subset at \(r\).  This proof is in \(\mathsf{ZF}\) and
requires no enumeration of a family of dense sets.
\end{proof}
\end{annexitem}
